\chapter{组合与风险：从协方差到成本感知组合}

\begin{itemize}
  \item \textbf{直接先修：}上册第 6、10、13、14、16、17 章；路线诊断见第 \ref{route:portfolio-risk} 节。
  \item \textbf{学习产出：}从协方差估计走到受约束优化、成本台账、VaR/ES 与压力测试，并交付一个可复现的两资产研究包。
  \item \textbf{研究边界：}合成小矩阵用于手算与失败诊断，不代表真实收益分布、可交易性或生产级优化。
\end{itemize}

\section{协方差、维度与因子分解}

令 $r\in\mathbb R^n$ 为同一时点的资产收益向量，$\mu=\mathbb E r$，协方差矩阵
\[
 \Sigma=\mathbb E[(r-\mu)(r-\mu)^\top].
\]
它必须对称半正定，因为任意权重 $w$ 都满足 $w^\top\Sigma w=\operatorname{Var}(w^\top r)\ge0$。样本协方差还隐含时间对齐、缺失值处理、收益口径和估计窗口；只检查数组形状不够。

因子模型写成 $r=Bf+\varepsilon$，若 $\operatorname{Cov}(f,\varepsilon)=0$，则
\[
 \Sigma=B\Sigma_f B^\top+D.
\]
本章的两资产 oracle 分别直接给出
\[
 \Sigma=\begin{pmatrix}0.04&0.01\\0.01&0.09\end{pmatrix}
\]
和独立的 $B=(1,0.25)^\top$、$\Sigma_f=(0.04)$、$D=\operatorname{diag}(0,0.0875)$；两条路径必须得到同一矩阵。若最小特征值为负，程序以“非半正定协方差”拒绝，而不是把错误矩阵继续送入优化器。

当资产数与样本量同阶甚至更大时，样本协方差可能奇异或病态。收缩估计
\[
 \widehat\Sigma_\lambda=(1-\lambda)\widehat\Sigma+\lambda T
\]
牺牲部分无偏性换取更小估计方差与更稳定的逆。目标矩阵 $T$ 和强度 $\lambda$ 必须由声明协议选择，不能事后挑到最好回测。

\begin{diagnostic}
“矩阵可逆”不等于“估计可靠”。应同时报告样本窗口、维度、最小特征值、条件数以及权重对输入扰动的敏感性。
\end{diagnostic}

\subsection{从收益表到可审计的协方差}

给定 $T$ 行、$n$ 列的收益矩阵，常用无偏样本协方差为
\[
 \widehat\Sigma=\frac{1}{T-1}
 \sum_{t=1}^{T}(r_t-\bar r)(r_t-\bar r)^\top.
\]
分母是 $T-1$ 还是 $T$、收益是简单收益还是对数收益、缺失日采用交集样本还是成对删除，都会改变矩阵。成对删除尤其危险：每个元素可能来自不同日期集合，拼出的矩阵不一定半正定。可审计实现应先冻结资产顺序和共同时间索引，再计算完整矩阵，并把对称化和特征值检查当成验收，而不是事后修饰。

本章的矩阵特征值约为 $0.0381$ 与 $0.0919$，条件数不高；它刻意让手算清楚。研究数据中若最小特征值接近机器精度，求逆会放大输入误差。此时应先追查数据与维度，再比较对角加载、收缩或因子结构，不能把 \verb|pinv| 当作自动修复。

Python 的 \verb|numpy.cov| 能计算矩阵，却不知道时点对齐和资产池规则。数学对象只在数据协议确定后才有明确含义；这正是从研究 notebook 迁移到可复现流程时最容易遗漏的一层。

\section{最小方差、约束与稳健性}

长仓、满仓的最小方差问题为
\[
 \min_w\ \frac12 w^\top\Sigma w,
 \qquad \mathbf1^\top w=1,\quad w\ge0.
\]
二资产内点解可直接手算。令 $w=(x,1-x)$，一阶条件给出
\[
 x^*=\frac{\sigma_2^2-\sigma_{12}}
 {\sigma_1^2+\sigma_2^2-2\sigma_{12}}.
\]
代入本章矩阵得到 $(8/11,3/11)=(0.727273,0.272727)$。这个分数是独立 oracle，不由优化实现反算。

加入边界约束时，KKT 乘子解释放宽约束的边际价值；乘子为零表示该约束在当前解不活跃。为真正观察不稳定性，本章另设病态矩阵
\[
 \Sigma_0=\begin{pmatrix}0.04&0.0399\\0.0399&0.04\end{pmatrix}.
\]
其最小方差第一资产权重为 $0.5$；只把第一项方差从 $0.04$ 改为 $0.0402$，权重就跳到 $0.25$。对扰动矩阵加入 $0.0075I$ 后，权重回到 $0.493506$。这组“基准--小扰动--正则化”对照才是稳健处理的执行证据；它不宣称 ridge 在所有数据上都更好。

为看清约束，写出带等式乘子 $\nu$ 与长仓乘子 $\lambda\ge0$ 的 Lagrangian：
\[
 \mathcal L(w,\nu,\lambda)
 =\tfrac12w^\top\Sigma w+\nu(\mathbf1^\top w-1)-\lambda^\top w.
\]
KKT 条件包括
\[
 \Sigma w+\nu\mathbf1-\lambda=0,\quad
 \mathbf1^\top w=1,\quad w\ge0,\quad\lambda\ge0,\quad
 \lambda_iw_i=0.
\]
最后一条互补松弛区分活跃与非活跃边界。若某资产权重为零，其乘子可能为正；这不是“优化器漏选了资产”，而是目标和约束共同给出的边界解。

稳健处理需要先声明不确定集。例如把期望收益写成 $\mu=\widehat\mu+u$，并限制 $u$ 落在椭球中，最坏情形优化会对暴露于高不确定方向的权重施加额外惩罚。它与 ridge 有相似的稳定化效果，但经济解释不同。教材实验只演示协方差对角加载，不把二者混为同一种方法。

均值--方差、最小方差和风险平价回答不同问题：前者需要脆弱的期望收益输入，第二个只使用风险输入，第三个平衡风险贡献而非直接最小化方差。Black--Litterman 等方法也没有消灭假设，只是把先验、观点和置信度显式化。

\section{成本感知再平衡与实施约束}

从当前权重 $w^-=(0.4,0.6)$ 再平衡到 $w$，本章使用
\[
 \max_w\quad \mu^\top w-\frac\gamma2w^\top\Sigma w
 -c\lVert w-w^-\rVert_1,
\]
并约束 $\mathbf1^\top w=1$、$0\le w_i\le0.6$。双边换手定义为 $\lVert w-w^-\rVert_1$；若机构使用单边口径，必须除以 2 并明确写出。

为保留可枚举 oracle，程序只在 0.1 网格上搜索。给定 $\mu=(0.08,0.04)$、$\gamma=1$ 和 $c=0.001$，最优可行权重为 $(0.6,0.4)$；双边换手为 $0.4$。资金规模十万元时，线性成本台账为
\[
 100000\times0.001\times0.4=40.
\]
这个账本独立于目标函数实现。

网格枚举还有一个教学价值：每个候选权重都能重算预期收益、方差和成本，因而可以定位第一处分歧。连续优化器给出相同答案后，枚举表仍应保留为小问题 oracle。若改变成本率或最大权重，答案应随输入变化；冻结某个结果而不绑定参数，会让“验证”失去意义。

成本的单位必须一致。这里 $c$ 是每单位双边换手的比例成本，目标中的成本是收益率口径，现金台账再乘资本。若把基点、百分数和小数混用，优化目标可能仍能运行，却会差一百倍或一万倍。研究报告应同时列比例成本与现金成本，使量纲错误更容易被发现。

可交易性是硬约束，不是成本项。停牌资产、借券不可得资产或涨跌停无法成交时，其权重不能随意改变。本章的负例尝试改变不可交易资产权重，必须以稳定诊断拒绝。真实系统还要处理整数股数、现金、最小订单、价差、冲击与部分成交。

\begin{note}[研究解释]
若忽略约束后再“把权重裁剪一下”，通常会同时破坏满仓、风险和最优性。应在同一优化问题中定义可行域，并在结果中报告活跃约束。
\end{note}

\section{VaR、ES 与压力测试}

令损失 $L=-R_p$。置信水平 $\alpha$ 下的风险价值为损失分布的分位数
\[
 \operatorname{VaR}_\alpha(L)=\inf\{\ell:\mathbb P(L\le\ell)\ge\alpha\}.
\]
为正确处理离散分布在分位点的概率质量，本章采用积分分位数定义
\[
 \operatorname{ES}_\alpha(L)=
 \frac{1}{1-\alpha}\int_\alpha^1\operatorname{VaR}_u(L)\,\mathrm du.
\]
对损失 $(-1,0,2,4)$ 和 $\alpha=0.75$，VaR 为 2，ES 为 4。若分位点存在重复值，算法会按尾部概率质量对边界观察加权，而不是按数组位置随意纳入或排除。

VaR 一般不是相干风险度量：它可能不满足次可加性；ES 在常见可积条件下具有相干性，但估计更依赖尾部样本。只有三个观察时，本章拒绝报告 95\% 尾部统计，而不是输出一个看似精确的数字。

离散样本的分位点可能有概率质量，因此软件包对 quantile 的插值默认值会产生不同 VaR。研究者必须固定“左分位数、线性插值或历史次序统计量”之一，并让手算例与实现使用同一口径。本章选择左分位数与积分分位数 ES；监管或机构定义可能采用不同离散估计约定。

VaR 的非次可加性可由稀有违约组合直观看出：每个头寸单独在给定置信水平下都可能显示 VaR 为零，但合并后“至少一个违约”的概率越过阈值，组合 VaR 反而为正。ES 关注尾部严重程度，较适合聚合风险，但并不自动解决模型错设、流动性枯竭或估计误差。

压力测试不从历史分位数机械推出。本章设两个资产的冲击为 $(-10,-2.5)$，对权重 $(0.6,0.4)$ 的损失为
\[
 -[0.6(-10)+0.4(-2.5)]=7.
\]
该情景是条件性问题，不是概率预测。报告应说明冲击来源、同时发生假设、二阶效应、流动性反馈和无法覆盖的情景。

轻尾正态模型可能低估跳跃、波动聚集和相关性上升。至少应并列历史/参数方法、压力情景和参数敏感性，而不是把一个 VaR 数字当作完整风险管理。

压力情景设计可以按三层组织：历史重演保留曾发生的联合冲击；假设情景改变利率、波动率、价差或相关性；反向压力测试则寻找使资本或流动性门槛失守的最小冲击。每层都要说明持仓是否静态、价格是否线性重估、相关性是否随压力变化，以及能否在压力中成交。

\begin{implementationnote}[从症状选择工具]
权重跳变先查输入条件数与约束活跃集；风险数字跳变先查损失符号、分位数口径和样本窗口；理论风险正常但现金亏损异常，则回到成本、成交和压力重估台账。不要用同一个“风险太大”标签覆盖三类原因。
\end{implementationnote}

\section{Route Capstone：成本感知组合与风险报告}

Capstone 从冻结 fixture 读取直接协方差和独立因子分解，依次核对最小方差、ridge 敏感性、受约束再平衡、成本台账、VaR/ES 与压力损失。三个负例分别守住协方差合法性、可交易性和尾部样本量。

冻结摘要见 \path{reports/lower-ch05-summary.md}。完整研究包必须报告输入哈希、估计协议、约束、换手口径、成本、风险定义、压力情景与不可声称的生产结论。

\begin{implementationnote}
\begin{lstlisting}[language=bash]
uv run python notebooks/lower/ch05_portfolio_risk.py \
  evidence/lower-ch05/oracle.json
\end{lstlisting}
\end{implementationnote}

\section{四级问题}
\begin{enumerate}[leftmargin=*]
  \item \textbf{口述概念：}为什么协方差矩阵必须半正定？可逆为何仍不等于可靠？
  \item \textbf{口述概念：}最小方差、均值--方差与风险平价分别优化什么？
  \item \textbf{笔试推导：}推导二资产长仓满仓最小方差权重。
  \item \textbf{笔试推导：}写出含 $L^1$ 换手成本和最大权重的再平衡问题，并解释活跃约束。
  \item \textbf{数值编程：}改变 ridge 强度，画出权重和组合方差的敏感性曲线。
  \item \textbf{数值编程：}用枚举核对本章再平衡台账，并分别报告单边、双边换手。
  \item \textbf{研究判断：}为何历史 VaR、参数 VaR 和压力测试不能互相替代？
  \item \textbf{研究判断：}设计停牌或借券不可得的故障注入，并说明恢复协议。
\end{enumerate}

